\newcommand{\DoNotLoadEpstopdf}{}
\documentclass[11pt,reqno]{amsart}
\usepackage[T1]{fontenc}
\usepackage{lmodern,amsmath,amssymb,amsthm,mathtools,mathrsfs,microtype}
\usepackage[margin=1.1in]{geometry}
\usepackage{xr-hyper}
\usepackage[hidelinks]{hyperref}
\externaldocument[supp-]{build/supplement}[supplement.pdf]
\hypersetup{pdftitle={Collision records, arithmetic, and raw local inversion in a triangular Lorentz gas},pdfauthor={Qian Qi}}
\setlength{\emergencystretch}{2em}
\allowdisplaybreaks[2]
\numberwithin{equation}{section}
\newtheorem{theorem}{Theorem}[section]
\newtheorem{lemma}[theorem]{Lemma}
\newtheorem{proposition}[theorem]{Proposition}
\newtheorem{corollary}[theorem]{Corollary}
\theoremstyle{definition}
\newtheorem{definition}[theorem]{Definition}
\newtheorem{inputstatement}[theorem]{Input statement}
\theoremstyle{remark}\newtheorem{remark}[theorem]{Remark}
\newcommand{\R}{\mathbb R}\newcommand{\Z}{\mathbb Z}
\newcommand{\one}{\mathbf1}\newcommand{\dd}{\,\mathrm d}
\newcommand{\esssup}{\operatorname*{ess\,sup}}
\newcommand{\diver}{\operatorname{div}}
\newcommand{\TV}{\operatorname{TV}}
\title[Collision records and raw local inversion]{Collision records, arithmetic, and raw local inversion in a triangular Lorentz gas}
\author{Qian Qi}
\date{October 12, 2026. A2-DYN, revision 84}
\subjclass[2020]{37D50, 37A50, 60F05, 28A75}
\keywords{Sinai billiards, exact collision records, local limits, backward residence, paired boundary traces, microscopic concentration, exact-label coarea transport}
\begin{document}
\raggedbottom
\begin{abstract}
We study the original exact-label pointwise local limit problem for
actual returns in a triangular Lorentz gas. The backward lifetime of
a one-sided outgoing face gives a residence kernel for its complete
first-exit assignment. This kernel converts the established local
surface-mass theorem into an essential-height bound of order
$r^{1/32}$ for all regular segments of length greater than $2r$.
For shorter segments away from a critical outgoing gradient, the
physical angular density transfers between the two true ends with
an error of order $r$ in the same collision-normalized essential norm.
Grazing incoming ends contribute no unpaid boundary trace. The
remaining short critical surface has a sharper local-mass bound.
We also identify a canonical microscopic concentration charge and
prove dyadic outer-scale rigidity of the fractional tail. The exact
labels, original section, guard, arithmetic reference and first-exit
selection remain unchanged. The final budget isolates paired short
nongrazing boundaries, short critical segments and a microscopic
charge; their remaining essential-height estimates are stated
separately from the new proved bounds.
\end{abstract}
\maketitle

\section{The unchanged physical problem}
Put $R_-=9/20$ and $R_+=47/100$, and let $R\in I=[R_-,R_+]$. Circular scatterers have centers in
$\Lambda=\Z(1,0)+\Z(1/2,\sqrt3/2)$. Write $T_R$ for the regular
collision map and $(\alpha,p)\in(\R/2\pi\Z)\times(-1,1)$ for its
coordinates; $p$ is the outgoing tangential velocity and
$\dd\nu=\dd\alpha\,\dd p/(4\pi)$. Physical flights satisfy
$\tau_-\le\tau_R\le\tau_+$, where $\tau_-=3/50$ and $\tau_+$
is the fixed finite-horizon bound.

The section is specified here to fix every event in the argument.
Put $h_0=1/200$. Let $Y_R$ be the union of the four open coordinate
squares of half-side $h_0$ centered at
\[
 (0,0),\quad(\pi/6,-1/2),\quad
 (-\theta_R,\sin\theta_R),\quad
 (\pi/3-\theta_R,\sin\theta_R).
\]
Here $\theta_R=2\arctan q_R$, and $q_R$ is the unique root in
$(49/100,53/100)$ of
\[
 8Rq(1+q^2)-4q(1-q^2)+\sqrt3(1-6q^2+q^4)=0.
\]
Set
\[
 U=\{|\alpha-\pi/6|<3/50,\ |p|<3/20\},\qquad
 Y_R^*=Y_R\cup U,\qquad c_*=\frac{91}{10000\pi}.
\]
The rectangles are disjoint and $\nu(Y_R^*)=c_*$. We always use
$\nu_R^*=c_*^{-1}\nu|_{Y_R^*}$, not a renormalization of a smaller
source.

At $m$ flights, let $L=L_{m,R}$ be their total length,
$K^{\rm c}_{m,R}$ their lattice displacement, and
$A_{m,R}=\sum_{j=0}^{m-1}\one_{Y_R^*}\circ T_R^j$. The event is
\begin{equation}\label{eq:v76-event}
 \Pi_{n,k,m,R}=Y_R^*\cap\{K^{\rm c}_{m,R}=k,\ A_{m,R}=n\}
                         \cap T_R^{-m}Y_R^*.
\end{equation}
The initial visit is counted; the terminal visit is tested separately.
On this event the actual $n$th return has collision count $m$ and
roof $L$. Denote its density by $p_{n,R}(k,m,t)$.

Write $c_j=\sqrt{1-p_j^2}$ and let $\Delta_j$ be the least clearance
of flight $j$ from any third disk, excluding its two incident disks.
Keep the original smooth function $0\le\Theta\le1$, zero on
$(-\infty,1]$ and one on $[2,\infty)$. With $q_0=47/53$,
$d_j=\min(j,m-j)$ and $d'_j=\min(j,m-1-j)$, set
\begin{align}
 H^\varepsilon&=
 \prod_{j=0}^{m}\Theta\!\left(\frac{c_j}{\varepsilon q_0^{d_j/4}}\right)
 \prod_{j=0}^{m-1}\Theta\!\left(\frac{\Delta_j}{\varepsilon q_0^{d'_j/4}}\right),
 \qquad W^\varepsilon=1-H^\varepsilon,\label{eq:v76-guard}\\
 b^\varepsilon(t)\dd t
   &=L_\#(W^\varepsilon\one_{\Pi_{n,k,m,R}}\nu_R^*).
 \label{eq:v76-source}
\end{align}
In particular first incidence and first clearance outside any selected
critical neighborhoods are still present. No section-distance factor
has been inserted into this guard.

\subsection{The arithmetic reference and the target}
For $y=(k_1,k_2,t,n)$ put
\[
 \mu_R=(0,0,\bar\tau_R,c_*),\quad
 Z=(y-m\mu_R)/\sqrt m,\quad
 \mathcal C_{m,R,n,k}(M)=\{t:|Z|\le M\}.
\]
We retain the finite spectral-peak reference, including peaks close to
but not on the unit circle. To specify its meaning, the collision
record observable is $(\kappa_1,\kappa_2,\tau,\one_{Y_R^*})$.
Let $\lambda_j(R,\vartheta)$ and $\mathsf P_j(R,\vartheta)$ be its
continued simple twisted eigenvalue and projection in each of the
fixed peak charts, and let $\chi_j$ be their frequency partition.
At the real maximum $a_j$ of $\log|\lambda_j|$ put
\[
 z_j=\lambda_j(R,a_j),\quad
 \mu_j=i^{-1}\nabla\log\lambda_j(R,a_j),\quad
 H_j=-\nabla^2\log\lambda_j(R,a_j),
\]
with $\operatorname{Re}H_j\ge aI_4$ for a common $a>0$. If
$\eta_R=\one_{Y_R^*}$ and $\ell$ denotes mass evaluation, put
$B_j=\chi_j(R,a_j)\ell(M_{\eta_R}\mathsf P_j(R,a_j)(\eta_R\nu))$.
Define
\begin{align}
 G_H(z)&=(2\pi)^{-4}\int_{\R^4}e^{-iv\cdot z-v^{\mathsf T}Hv/2}\dd v,
 \notag\\
 \mathcal L_{m,R}(y)&=\operatorname{Re}\left[
 c_*^{-1}\sum_j e^{-ia_j\cdot y}z_j^m B_j
      G_{H_j}\left((y-m\mu_j)/\sqrt m\right)\right].
 \label{eq:v76-reference}
\end{align}
Only the charts active at $R$ are summed. This is the same transition
reference as in the retained spectral theorem. The existence and
uniform estimates for these scalar data belong to that analytic
input, not to the transport argument below. No frequency branch or
arithmetic zero class is deleted.

The governing assertion remains
\begin{equation}\label{eq:v76-target}
 \sup_{R,\,1\le n\le m,\,k\in\Z^2}
 \esssup_{t\in\mathcal C_{m,R,n,k}(M)}
 |m^2p_{n,R}(k,m,t)-\mathcal L_{m,R}(k_1,k_2,t,n)|\longrightarrow0.
\end{equation}
Let $\mathcal E_M$ denote this error's collision limsup. An empty
central target contributes zero. The order of the uniformity is
$\sup_R\esssup_t$, so exceptional density sets may depend on $R$.

\subsection{The complete surface-height estimates proved here}
For each one-sided nongrazing outgoing face, let $l$ be its backward
lifetime in the original regular exact event. Write
$\overline{\mathcal F}_r$ for the complete first-exit surface assignment
averaged over roof increments in $(0,r)$. The source is not replaced
by a boundary collar. Its restriction to $l>2r$ is denoted by
$\mathcal L_r$. Among the remaining segments, let
$\mathcal S_r^{\rm nc}$ have outgoing $|g|^2\ge8\Lambda_\circ r$,
and let $\mathcal J_r$ carry their incoming physical traces, paired
with the same outgoing faces and the same exit-window kernel.

\begin{theorem}[Residence and two-ended optical transfer]
\label{thm:v84-main}
Under the explicit inherited inputs below, for every sufficiently
small previously fixed $r>0$,
\begin{align}
 \limsup_m\sup_{R,n,k}m^2\|\mathcal L_r\|_\infty
                          &\le Cr^{1/32},\label{eq:v84-main-long}\\
 \limsup_m\sup_{R,n,k}m^2
       \|\mathcal S_r^{\rm nc}-\mathcal J_r\|_\infty
                          &\le Cr.\label{eq:v84-main-pair}
\end{align}
The complete short noncritical part with grazing incoming end also
has height at most $Cr$. All comparisons retain the original exact
label and every incident physical branch; no word-count multiplier
is used.
\end{theorem}
The first conclusion is Theorem~\ref{thm:v84-long-height}. It uses a
persistence property of the actual surface kernel, not an inference
from small mass alone. The second is
Theorem~\ref{thm:v84-paired-transfer}; it estimates the change of
physical angular density between the two ends, including both angular
Jacobian factors. The positive lower gradient used inside that proof
is a function of the fixed roof width only.

The preceding untruncated fractional transport and local surface-mass
results are retained with their proofs. The new argument uses them to
control the complementary long-residence source in essential height.
The still unresolved sources are short paired nongrazing boundaries,
short critical segments, and a canonical microscopic concentration
charge. Theorem~\ref{thm:v84-endpoint} gives their exact role in the
unchanged pointwise assertion \eqref{eq:v76-target}.

\section{Explicit analytic and geometric inputs}
We collect the indispensable inherited assertions in the exact form
used by the new proofs. Their source theorems and operator dependencies
are specified here and in Appendix~\ref{app:v83-kernels}. The new
potential and first-exit estimates do not assume a surface spectral
insertion theorem or an unweighted pointwise local bound.

\begin{inputstatement}[Physical flow and source comparison]
\label{input:v83-geometric}
On the regular part $\Omega$ of each exact event, endpoint angles
$q=(\alpha_0,\alpha_m)$ give $F=L$, $g=\nabla_qF=R(-p_0,p_m)$,
$K=D_q^2F$ and $\vartheta=\arg g$. Put
\[
 \lambda_\circ=9\sqrt{19}/200,\qquad \Lambda_\circ=47/6,\qquad
 \omega=|\partial_{p_0}\alpha_m|^{-1},\qquad
 \chi=|\det D_q(F,\vartheta)|.
\]
The bounds $\lambda_\circ I\le K\le\Lambda_\circ I$ and
$\lambda_\circ\le\chi\le\Lambda_\circ$ hold uniformly.
The explicit field
\[
 X_q=\frac{K^{-1}g}{g^{\mathsf T}K^{-1}g}
\]
is stopped at true physical and section-decision boundaries, not at
chart seams. Its flow satisfies $XF=1$, $X\vartheta=0$ and
\begin{equation}\label{eq:v83-radial-input}
 2\lambda_\circ s\le |g(\Phi_sx)|^2-|g(x)|^2
                                      \le2\Lambda_\circ s.
\end{equation}
Its endpoint-curve length is at most $C\sqrt s$. For
$I_m=1+\sum_{i=1}^{m-1}q_0^{d_i}/c_i$ one has
$|X\omega|\le C\omega I_m/|g|$ on regular survivors. At a genuine
grazing exit some internal incidence tends to zero and $\omega$ tends
to zero. Forward interior critical termination cannot occur.

Every starting point with forward lifetime at most $r$ belongs to a
layer of one of the following forms, at an actual collision or flight:
\begin{equation}\label{eq:v83-layer-input}
 c_j\le C_Eq_0^{d_j}\sqrt r,\quad
 \Delta_j\le C_Eq_0^{d'_j}\sqrt r,\quad
 \operatorname{dist}(T_R^jx,\partial Y_R^*)\le C_Eq_0^{d_j}\sqrt r.
\end{equation}
The physical coarea covariance, proved in the appendix in the convention
used here, has Jacobian $J_h=(\omega(\Phi_hx)/\omega(x))\chi(x)/\chi(\Phi_hx)$.
Writing $E_h^\varepsilon$ for all weighted starting exit mass and
\[
 O_h^\varepsilon(t)=\int_{S_h(t)} W^\varepsilon(x)
            [1-\omega(\Phi_hx)/\omega(x)]_+\,d\mu_t(x),
\]
the complete source satisfies, for each fixed $r$ and almost every $t$,
\begin{equation}\label{eq:v83-shift-input}
 b^\varepsilon(t)\le2A_rb^\varepsilon(t)+\bar E_r^\varepsilon(t)
       +\bar O_r^\varepsilon(t)
       +C_\Theta(1+\varepsilon^{-1})\sqrt r\,A_rp(t),
 \qquad A_rf(t)=\frac1{2r}\int_{t-r}^{t+r}f(s)\,ds.
\end{equation}
\end{inputstatement}
These are the endpoint-field, all-contact motion, uniform cubic-action
and separated-volume results in retained modules 174, 177, 180--182.
The guard term in \eqref{eq:v83-shift-input} is the complete product
guard, not a protective subword source. The derivative and continuation
assertions are confined to regular physical trajectories.

\begin{inputstatement}[Complete local source estimates]
\label{input:v83-local}
For every fixed auxiliary band $B_*$, constants $C_{B_*}<\infty$ and
$\rho_{B_*}<1$ give the following bounds uniformly in radius, exact
labels, interval $J$ and mark $0\le j\le m$. Let $S_{R,s}$ be the
union of strips of width $Cs$ about the finitely many coordinate lines
supporting $\partial Y_R^*$. For $s$ small,
\begin{align}
 m^2\nu\{K_m^{\rm c}=k,A_m=n,L_m\in J,
                       T_R^jx\in S_{R,s}\}
 &\le Cs^{1/16}(|J|+B_*^{-1})
   +C_{B_*}(1+B_*|J|)m^2\rho_{B_*}^{m/2},\label{eq:v83-strip-input}\\
 m^2\int_Jb^\delta(t)\,dt
 &\le C\delta^{1/16}(|J|+B_*^{-1})
   +C_{B_*}(1+B_*|J|)m^3\rho_{B_*}^{m/2}.\label{eq:v83-guard-input}
\end{align}
The local extension to a fixed remaining compact range of $\delta$
is included. At auxiliary parameter one $W^1=1$, so $b^1=p$.

Put $z=|g|$, $\alpha=129/128$ and, for $0<\eta<1/2$, define the
actual nonnegative source measures
\begin{align}
 v^{\eta,b}(t)\,dt&=L_\#(I_m z^{-\eta}\one_{\{0<z<b\}}
                                       \one_\Pi\nu_R^*),\notag\\
 v^\eta(t)\,dt&=L_\#(I_m z^{-\eta}\one_\Pi\nu_R^*).
                                                   \label{eq:v83-fractional-input}
\end{align}
They are finite and absolutely continuous at every fixed count, and
for $0<b<b_1$,
\begin{equation}\label{eq:v83-fractional-finite-input}
 m^2\int_Jv^{\eta,b}(t)\,dt
 \le C_\eta b^{\alpha-\eta}(|J|+B_*^{-1})
  +C_{B_*,\eta}(1+B_*|J|)m^2b^{1/2-\eta}\rho_{B_*}^{m/4}.
\end{equation}
Adding the bounded-gradient part gives, for every fixed $L>0$,
\begin{equation}\label{eq:v83-fixed-width-input}
 \limsup_m\sup_{R,n,k}\sup_{|J|=L}m^2\int_Jv^\eta\le C_\eta L,
 \qquad
 \limsup_m\sup_{R,n,k}\sup_{|J|=L}m^2\int_Jv^{\eta,b}
                                      \le C_\eta b^{\alpha-\eta}L.
\end{equation}
Every auxiliary band is fixed before the collision limit and enlarged
only afterwards; leading constants are independent of that band and
of the previously fixed interval length.
\end{inputstatement}
Equation~\eqref{eq:v83-strip-input} is Supplement
Theorem~\ref{supp-thm:v48-marked-strip-local}. The physical source
bound \eqref{eq:v83-guard-input} is Supplement
Theorem~\ref{supp-thm:v49-physical-remainder-local}, in particular its
finite-window estimate. The fractional bounds are the positive
layer-cake conclusions of retained module 187, stated in full above.
The joint spectral and common-space assumptions used there are recorded
in the appendix. No estimate for an untruncated first inverse gradient
or for a singular roof convolution is part of this input.

\begin{inputstatement}[The scalar positive remainder]
\label{input:v83-scalar}
For large fixed $B$, with $\varepsilon(B)=A_0B^{-1/12}$,
\begin{equation}\label{eq:v76-remainder}
 m^2p-\mathcal L=m^2b^{\varepsilon(B)}+e_B,\qquad
 \limsup_m\sup_{R,n,k}\esssup_{\mathcal C(M)}|e_B|
                                      \le C_MB^{-1/192}.
\end{equation}
The source $b^\varepsilon$ is nonnegative and has a compactly supported
$BV$ density at fixed count. No collision-uniform variation bound is
assumed. Supplement Theorem~\ref{supp-thm:v51-positive-error} gives
this remainder with the same complete finite-peak reference.
\end{inputstatement}

\input{core/190_untruncated_optical_potentials}
\input{core/191_complete_first_exit_surface_mass}
\input{core/192_ordered_fractional_tail_and_endpoint}

\input{core/193_backward_residence_surface_height}
\input{core/194_paired_boundary_optical_transfer}
\input{core/195_canonical_concentration_and_endpoint}


\section{The original scalar objective and downstream normalization}
The governing assertion remains \eqref{eq:v76-target}. The long-residence
surface and the short noncritical optical transfer now have proved
essential-height bounds. Their complements are not removed: the
paired nongrazing trace and short critical surface remain inside the
same exact-label norm. The latter has the quantitative local-mass
bound \eqref{eq:v84-critical-mass}, which is not its height bound.
The canonical microscopic charge in \eqref{eq:v84-canonical-charge}
requires a separate vanishing estimate; dyadic outer-scale rigidity
shows precisely why reducing the reconstruction width cannot replace
that estimate.

The only active dyadic convention in this article is the offset formula
\eqref{eq:v83-dyadic-tail}. The retained module 189 is marked as
superseded at its point of inclusion in the supplement, and its original
bytes remain in the historical source. Its ambiguous integration
notation is not an active theorem used here.

After scalar closure, the positive path-remainder equality in
Supplement Theorem~\ref{supp-thm:v73-scalar-path-equality} supplies
the same-roof path numerator. Ratios are formed only on positive
arithmetic reference classes; the actual-return clock and forward
likelihood retain their separate inherited hypotheses. The unchanged
scalar-first order is not reversed by any surface comparison.

\appendix
\input{appendices/v83_common_kernels_and_dependencies}
\input{journal-references}
\end{document}
